\documentclass[11pt]{article}
\usepackage{amsmath,amssymb}
\usepackage[margin=1in]{geometry}
\usepackage{booktabs}
\usepackage{tabularx}
\usepackage{hyperref}
\setlength{\emergencystretch}{2em}

\title{Representation Accessibility Changes the Compute Frontier:\\
Toward General Scaling Laws for Structured Scientific Reasoning}
\author{ORION Paper IX Ultimate-Successor Working Manuscript}
\date{August 2026}

\begin{document}
\maketitle

\begin{abstract}
A model can receive the same semantic information through representations that differ dramatically in how easily the task-relevant relations can be accessed. ORION Paper IX targets the strongest upward claim: across task families, relational complexities, model families, sample sizes, and inference budgets, exposing the right task coordinates changes the model-capacity, data, and test-time-compute frontier required to reach a fixed verified quality target. Existing P9 evidence of information-equivalent relational-versus-flat separation, interaction-order effects, fresh controlled replication, capacity/sample substitution, semantic-orbit stability, access-degree recovery, inductive-bias substitution, and bounded held-out-domain transfer is preserved as the starting foundation. P9-U does not stop at another representation benchmark. It prospectively freezes a factorial law over relational complexity $k$, representation $R$, model scale $S$, sample budget $N$, and inference budget $C$, and estimates critical frontiers $S^*(k,q)$, $N^*(k,q)$, and $C^*(k,q)$ for verified quality $q$. Strong baselines include equal-token flat representations, classical/exact query interfaces, adaptive-compute-only policies, architecture priors, graph/relational models, and donor-complete representation systems. The intended terminal is a replicated cross-family scaling law showing that representation accessibility changes the growth rate or crossing point of required capacity/compute, not merely the mean score, on at least two open-weight model families and heterogeneous formal/procedural domains. Historical negative P9/P10 results remain immutable; the existing negative-to-positive successor programme is used to discover more general access coordinates and operator-learning mechanisms when a scale cell fails.
\end{abstract}

\section{Largest scientific question}
For task complexity $k$, representation $R$, model family $F$, model scale $S$, sample budget $N$, inference budget $C$, and quality target $q$, define the verified performance function
\[
Q(k,R,F,S,N,C).
\]
P9-U asks whether the representation changes the resource frontier itself:
\[
S^*(k,q\mid R),\quad N^*(k,q\mid R),\quad C^*(k,q\mid R),
\]
not just a single fixed-budget accuracy point.

The strongest claim is that information-equivalent representations can induce systematically different accessibility scaling laws, and that exposing the correct task coordinates can let materially smaller/cheaper systems reach targets that larger flat systems cannot reach at the same budget.

\section{Upward-generalization principle}
Graph neural networks, relational inductive bias, query interfaces, structured state, architecture priors, test-time scaling, and exact symbolic computation are donor mechanisms. P9-U admits all of them to the comparator. ORION earns superiority only by a general cross-regime law that predicts when representation changes the resource frontier and survives stronger models, more compute, multiple domains, and multiple model families.

The intended terminal is \textbf{GENERAL_REPRESENTATION_ACCESSIBILITY_SCALING_LAW}.

\section{Immutable predecessor evidence}
The current bounded P9 paper and successor studies remain immutable, including positive information-equivalent relational-access results and any negative/donor-subsumed cells. The direct Qwen scaling run under the active successor programme is new evidence only; it cannot retroactively rewrite the earlier paper.

Negative P9/P10 results are routed through the existing ORION negative-to-positive programme \#586. In particular, known-operator composition, failure-applicability, serialization/access geometry, and hypothesis-class inadequacy remain distinct failure classes rather than being pooled.

\section{Factorial scaling protocol}
Freeze a prospective design over:
\begin{enumerate}
    \item relational/task complexity $k$;
    \item representation family $R$;
    \item model family $F$;
    \item model scale $S$;
    \item training/sample budget $N$;
    \item inference/search budget $C$;
    \item verified quality target $q$.
\end{enumerate}

The primary analysis estimates on-grid threshold/crossover quantities rather than fitting a post-hoc power law. A scaling exponent is reported only if its functional form is frozen before protected outcomes and diagnostics support it.

\section{Representation families}
At minimum include:
\begin{enumerate}
    \item information-equivalent flat/text serialization;
    \item reversible indexed serialization;
    \item typed relation tuples/set representation;
    \item native typed graph/relational state;
    \item task-sufficient query-matched interface;
    \item architecture-prior baseline that receives structure through the model rather than the input;
    \item adaptive-compute-only arm;
    \item representation-plus-adaptive-compute arm.
\end{enumerate}

Equal-token/length controls, symbol/order reminting, semantic-orbit controls, and exact round-trip information checks are mandatory.

\section{Task families}
The study spans at least:
\begin{enumerate}
    \item formal or mathematical relational reasoning;
    \item algorithmic/program/procedural reasoning;
    \item graph or dependency reasoning;
    \item a non-formal procedural/scientific workflow domain;
    \item optional later scientific episodes only after exact/controlled mechanisms are stable.
\end{enumerate}
At least one domain is held out entirely from representation/mechanism tuning.

\section{Strong baselines}
The comparator set includes:
\begin{enumerate}
    \item larger flat/text model;
    \item equal-parameter structured model;
    \item graph/relational architecture;
    \item exact/classical task-sufficient query engine where applicable;
    \item adaptive inference/test-time compute;
    \item representation compiler whose preprocessing cost is fully charged;
    \item donor-complete system combining structure, architecture prior, and adaptive compute;
    \item second open-weight model family for replication.
\end{enumerate}

\section{Primary hypotheses}
\paragraph{H1, scale-frontier shift.}
For increasing relational complexity, the critical model scale required to reach quality $q$ grows more slowly under the task-accessible representation than under information-equivalent flat representation.

\paragraph{H2, inference-frontier shift.}
At matched model scale, the critical inference budget required to reach $q$ is lower under the task-accessible representation.

\paragraph{H3, data-frontier shift.}
At matched model and inference budget, the sample threshold required to reach $q$ is lower under the accessible representation.

\paragraph{H4, non-format explanation.}
The effect survives equal-length/token, reminting, semantic-orbit, and stronger-model attacks and is not explained by preprocessing computing the answer.

\paragraph{H5, cross-family replication.}
The frontier shift replicates prospectively on a second open-weight model family and at least two heterogeneous task domains.

\paragraph{H6, access-mechanic expansion.}
At least one negative/tied scale regime yields a newly learned access coordinate, local operator, or query interface that expands protected reach beyond the pre-freeze representation basis and transfers to fresh tasks.

\section{Primary outcomes}
Report:
\begin{enumerate}
    \item critical model scale $S^*(k,q)$;
    \item critical sample budget $N^*(k,q)$;
    \item critical inference budget $C^*(k,q)$;
    \item verified task quality at matched resources;
    \item representation/token/memory cost;
    \item compiler/preprocessing cost;
    \item robustness to reminting/permutation/orbit controls;
    \item cross-domain and cross-model-family effect;
    \item failure cells and no-crossing cells.
\end{enumerate}

\section{Adversarial controls}
Mandatory attacks include:
\begin{enumerate}
    \item structure only shortens the prompt;
    \item larger models become invariant and erase the gap;
    \item extra preprocessing computes the target;
    \item test-time compute alone closes the gap;
    \item architecture prior substitutes fully for representation;
    \item representation effect disappears after exact query matching;
    \item formatting priors explain the result;
    \item train/test templates leak relational structure;
    \item post-hoc chosen scale grid manufactures a crossing.
\end{enumerate}

\section{Negative-result recursion as upward access search}
Every failed scale cell is preserved. The programme distinguishes representation insufficiency, access-operator absence, architecture limitation, model-family prior, resource regime, task non-identifiability, and donor-complete exact solution. It then searches materially different successor mechanisms rather than tuning the same representation.

The existing #586 programme already owns this doctrine. Known-operator negatives may escalate to learned local operators; access-geometry failures may escalate to query-matched or representation-compilation studies; hypothesis-class failures may escalate only after lower-level access causes are ruled out. A new child issue is opened only for a materially distinct access/mechanism family.

\section{Claim ladder and current status}
The intended progression is: direct open-weight scaling result; second-family replication; frozen multi-factor frontier; cross-domain transfer; hostile explanation controls; \textbf{GENERAL_REPRESENTATION_ACCESSIBILITY_SCALING_LAW} with H6 access-mechanic expansion.

This TeX file is prospective and does not claim the scaling law before the frozen real-system study executes.

\end{document}
